\subsection{Clinical Implications}

From a clinical perspective, three key findings emerge regarding the scope of seizure coverage, the readiness of detection devices, and the current state of forecasting technology.

\subsubsection{Seizure Type Coverage Gap}

Current wearable approaches address only a subset of seizure types. Accelerometer-based systems detect motor manifestations but cannot identify seizures without movement components. The ECG-based approach by \textcite{reintjesECGBasedDetectionEpileptic2025} examined whether autonomic signals could detect non-motor seizures. However, the observed sensitivity-FPR tradeoff (2.44\% to 82.8\% sensitivity at FPR ranging from 0.11/h to 65.62/h) prevents clinical interpretation. The lowest false alarm rate detected fewer than 3\% of seizures, while configurations achieving acceptable sensitivity generated false alarm burdens orders of magnitude above clinically acceptable thresholds. This fundamental tradeoff demonstrates that single-lead ECG anomaly detection cannot currently provide clinically reliable seizure detection.

Forecasting approaches may theoretically address this gap by relying on autonomic modalities rather than movement sensors. However, the consistent AUC ceiling around 0.75 suggests fundamental limits in predicting seizures without EEG biomarkers. Approximately 30\% of people with epilepsy have seizures that would not be detected by current wearable approaches. This coverage gap represents a substantial unmet need that may require invasive monitoring or alternative biomarkers.

\subsubsection{Detection Readiness}

The detection literature demonstrates that clinically acceptable FPR is achievable with accelerometer-based approaches (see Results). Two studies meet the ILAE Phase 3 benchmark using ACC data, suggesting that motor seizure detection has reached sufficient maturity for clinical translation. However, this achievement remains concentrated in controlled EMU settings rather than real-world home environments.

The EMU-to-home translation gap represents a critical barrier. The two studies achieving benchmark FPR were validated in hospital settings where patient activities are restricted and controlled. Home environments introduce motion patterns from daily living that do not occur in EMUs, including exercise, household chores, and other routine activities. The substantial FPR increase observed in home validation (0.165/h versus 0.0054/h in EMU settings) suggests that algorithm robustness to real-world activities must become a research priority.

\subsubsection{Forecasting Status}

The forecasting literature shows a consistent performance ceiling around AUC 0.74 to 0.75 across different methods, modalities, and patient populations (see Results). This convergence suggests a fundamental limit on pre-ictal predictability using non-invasive wearable sensors. Unlike detection, where clear clinical benchmarks exist, forecasting lacks standardized validation guidelines. Without established performance targets, the clinical utility of AUC 0.75 remains uncertain.

The personalization strategy fundamentally affects forecasting utility. Global LOSO validation achieves only 43\% patient success, meaning more than half of patients receive no benefit from population-based forecasting. Patient-specific approaches achieve 100\% success but require weeks to months of individual data collection before the system becomes useful. This cold-start problem represents a substantial barrier to clinical adoption, as newly diagnosed patients cannot wait months for protection.

No device has regulatory approval specifically for seizure forecasting. Forecasting devices would need to establish clinical utility and safety through dedicated trials. The lack of clear benchmarks complicates pathway to regulatory approval.
